%==============================================================
% BAB III  METODE
%==============================================================

\chapter{METODE}

%--------------------------------------------------------------
\section{Data dan Strategi Pembagian Data}
%--------------------------------------------------------------

\subsection{Dataset Utama: Zhao2024}

Dataset utama yang digunakan dalam penelitian ini adalah koleksi citra fundus bayi yang
diterbitkan oleh \citet{Zhao2024}, selanjutnya disebut Zhao2024. Dataset ini terdiri dari
1.099 citra fundus bayi prematur dalam lima kelas: Normal, Stage 1, Stage 2, Stage 3, dan
bekas luka laser (\textit{laser scars}). Distribusi jumlah citra untuk setiap kategori disajikan pada Gambar \ref{fig:class_distribution}. 

Untuk memberikan pemahaman visual mengenai karakteristik klinis pada setiap kelas ROP dalam dataset Zhao2024, contoh citra untuk masing-masing kategori ditunjukkan pada Gambar \ref{fig:dataset_samples}. Gambar tersebut memperlihatkan struktur retina yang normal hingga kelainan patologis seperti garis batas (\textit{demarcation line}) tipis pada Stage 1, tonjolan jaringan (\textit{ridge}) yang menebal pada Stage 2, pertumbuhan neovaskularisasi ekstraretina ke arah \textit{vitreous} pada Stage 3, dan area bekas tindakan medis fotokoagulasi laser pada kelas \textit{laser scars}.

\begin{figure}[H]
  \centering
  \includegraphics[width=0.7\linewidth]{images/class_distribution}
  \caption{Distribusi jumlah citra untuk setiap kategori pada dataset utama}
  \label{fig:class_distribution}
\end{figure}

\begin{figure}[H]
  \centering
  \includegraphics[width=\linewidth]{images/dataset_samples}
  \caption{Contoh citra fundus mata bayi prematur untuk masing-masing kelas dalam dataset Zhao2024}
  \label{fig:dataset_samples}
\end{figure}

Analisis mendalam terhadap dataset menemukan adanya 12 citra yang identik secara byte
(\textit{byte-identical duplicates}) dalam 6 kelompok duplikat. Yang lebih penting, tiga
dari kelompok tersebut merupakan konflik label antar-kelas: citra yang persis sama tercatat
di bawah dua kelas yang berbeda. Pasangan yang teridentifikasi adalah sebagai berikut:
\begin{center}
  \textit{Stage\_2\_ROP\_154.jpg} $\equiv$ \textit{Stage\_3\_ROP\_210.jpg} \\
  \textit{Stage\_2\_ROP\_52.jpg} $\equiv$ \textit{Stage\_3\_ROP\_254.jpg} \\
  \textit{Stage\_1\_ROP\_42.jpg} $\equiv$ \textit{Stage\_2\_ROP\_4.jpg}
\end{center}
Duplikat ini dibiarkan ada dalam dataset dan ditangani melalui protokol pengelompokan (\textit{grouping}), yaitu seluruh berkas duplikat ditempatkan ke dalam kelompok yang sama sehingga tidak pernah terbagi antara himpunan pelatihan dan pengujian.

\subsection{Strategi Pembagian Data dan Pengelompokan}

Karena dataset Zhao2024 tidak menyediakan identitas pasien pada nama file, pembagian data
secara langsung berbasis identitas pasien tidak dapat dilakukan. Sebagai solusi, analisis
kemiripan antar citra menggunakan korelasi silang ternormalisasi (\textit{Normalized Cross-
Correlation}, NCC $\geq$ 0,99) dilakukan untuk mendeteksi pasangan citra yang berasal dari
mata yang sama. Proses ini menghasilkan 721 kelompok dari total 1.099 citra yang kemudian
disimpan dalam berkas \textit{clean\_manifest.json}.

Validasi silang 5-\textit{fold} menggunakan \textit{StratifiedGroupKFold}---varian validasi silang yang menjamin distribusi kelas tetap merata di setiap \textit{fold} sekaligus memastikan seluruh anggota satu kelompok tidak terbagi antara himpunan pelatihan dan pengujian---diterapkan agar setiap kelompok selalu berada sepenuhnya di satu sisi batas pelatihan--pengujian. Strategi ini memastikan bahwa model benar-benar mempelajari pola penyakit, bukan pola citra dari mata tertentu.

\begin{figure}[H]
  \centering
  \includegraphics[width=0.85\linewidth]{images/group_fold_distribution}
  \caption{Distribusi kelas pada setiap \textit{fold} setelah pembagian \textit{group-aware}}
  \label{fig:group_fold_distribution}
\end{figure}

%--------------------------------------------------------------
\section{Pembuatan Representasi Softmap Spasial}
%--------------------------------------------------------------

Representasi \textit{softmap} spasial dibuat untuk memberikan tiga jenis bukti visual kepada
model klasifikasi: tekstur pembuluh darah, struktur area tepi atau \textit{ridge}, dan konteks
intensitas kanal hijau. Seluruh kanal diproses pada citra yang lebih dulu diperkecil dengan
batas sisi terpanjang 768 piksel, kemudian diubah ke ukuran akhir $224 \times 224$ piksel
sebagai masukan \textit{TinyResNet}.

\subsection{Estimasi FOV dan Normalisasi Kanal Hijau}

Area \textit{field of view} (FOV) diperkirakan dari citra abu-abu dengan ambang intensitas
rendah, dilanjutkan operasi \textit{morphological closing}, pengisian lubang, dan pemilihan
komponen terhubung terbesar. Masker FOV ini digunakan untuk meniadakan latar luar retina
pada seluruh kanal. Kanal hijau dipilih sebagai kanal dasar karena kontras antara pembuluh
darah dan jaringan retina paling jelas pada kanal ini. Setelah FOV diterapkan, kanal hijau
ditingkatkan menggunakan CLAHE dan dinormalisasi persentil ke rentang $[0, 1]$.

\subsection{Softmap Pembuluh Darah Berbasis Plain Gabor}

Kanal pembuluh darah dibuat dengan menonjolkan struktur tubular gelap pada kanal hijau.
Pertama, kanal hijau yang telah diberi CLAHE dibalik intensitasnya sehingga pembuluh darah
gelap menjadi respons terang. Selanjutnya, bank filter Gabor multiskala dan multiorientasi
diterapkan pada citra tersebut. Skala yang digunakan adalah $\sigma = 1{,}5; 2{,}5; 3{,}5;
5{,}0$ piksel dengan panjang gelombang 3, 5, 7, dan 10 piksel. Orientasi filter diambil dari
$0^\circ$ hingga $165^\circ$ dengan langkah $15^\circ$. Respons maksimum dari seluruh filter
digunakan sebagai peta pembuluh darah awal, kemudian diperhalus dengan median blur $7 \times
7$, diberi CLAHE kedua, dan dinormalisasi kembali ke rentang $[0, 1]$.

\begin{figure}[H]
  \centering
  \includegraphics[width=\linewidth]{images/vessel_preprocessing_steps}
  \caption{Tahapan pembentukan \textit{softmap} pembuluh darah berbasis plain Gabor}
  \label{fig:vessel_preprocessing_steps}
\end{figure}

\subsection{Softmap Ridge Berbasis Meijering dan Oriented Top-Hat}

Kanal \textit{ridge} dirancang untuk memperkuat struktur garis batas, tonjolan, dan area tepi
yang relevan secara klinis pada Stage 1 hingga Stage 3. Respons pertama diperoleh dari filter
Meijering berbasis Hessian pada kanal hijau CLAHE dengan skala $3, 5, 7, 9,$ dan $11$ piksel.
Respons kedua diperoleh dari \textit{oriented top-hat} dengan elemen struktur garis pada 12
orientasi. Kedua respons digabungkan sebagai:
\begin{equation}
  \text{ridge\_softmap} =
  \operatorname{norm}\left(0{,}60 \times \text{meijering\_ridge}
  + 0{,}40 \times \text{oriented\_tophat}\right).
  \label{eq:ridge_softmap}
\end{equation}

\begin{figure}[H]
  \centering
  \includegraphics[width=\linewidth]{images/ridge_preprocessing_steps}
  \caption{Tahapan pembentukan \textit{softmap ridge} berbasis Meijering dan \textit{oriented top-hat}}
  \label{fig:ridge_preprocessing_steps}
\end{figure}

\subsection{Komposit Softmap Tiga Kanal}

Masukan \textit{softmap} akhir disusun dari tiga kanal kontinu:
\begin{equation}
  X_{\text{softmap}} =
  [\text{plain\_gabor\_vessel},\ \text{ridge\_softmap},\ \text{clahe\_green}].
  \label{eq:softmap_input}
\end{equation}

Ketiga kanal tersebut dimasker oleh FOV dan dinormalisasi sebelum diubah ukurannya menjadi
$224 \times 224$ piksel. Dengan susunan ini, model menerima informasi pembuluh darah,
informasi \textit{ridge}, dan konteks intensitas retina secara terpisah tetapi tetap spasial.

\subsection{Evaluasi Pendukung Segmentasi Struktural}

Selain digunakan sebagai representasi kontinu untuk klasifikasi, kanal struktural juga
dievaluasi secara kuantitatif terhadap anotasi segmentasi dari dataset Agrawal2021. Kanal
pembuluh darah dievaluasi terhadap HVDROPDB-BV, sedangkan kanal \textit{ridge} dievaluasi
terhadap HVDROPDB-RIDGE. Evaluasi ini digunakan sebagai pemeriksaan pendukung bahwa peta
struktur yang dihasilkan masih mengikuti anotasi spasial retina. Untuk evaluasi biner, peta
respons terlebih dahulu dibinarisasi, kemudian dibandingkan dengan \textit{ground truth}
menggunakan \textit{Dice coefficient}:
\begin{equation}
  \text{Dice} = \frac{2 \times |A \cap B|}{|A| + |B|},
  \label{eq:dice}
\end{equation}
dengan $A$ menyatakan masker prediksi dan $B$ menyatakan masker anotasi pakar.

%--------------------------------------------------------------
\section{Representasi Citra Input}
%--------------------------------------------------------------

Penelitian ini membandingkan dua representasi citra input untuk menguji pengaruh
pemrosesan citra spasial terhadap performa model klasifikasi. Gambar \ref{fig:pipeline_diagram} menunjukkan diagram alir \textit{pipeline} penelitian.

\begin{figure}[H]
  \centering
  \includegraphics[width=\linewidth]{images/pipeline_diagram}
  \caption{Diagram alir \textit{pipeline} penelitian dari representasi input hingga klasifikasi}
  \label{fig:pipeline_diagram}
\end{figure}

\vspace{2mm}
\noindent\textbf{Citra RGB Mentah (\textit{Raw RGB})} \\
Citra fundus asli yang hanya diubah ukurannya menjadi $224 \times 224$ piksel dan dinormalisasi nilai pikselnya ke rentang $[0, 1]$. Representasi ini berfungsi sebagai pembanding dasar (\textit{baseline}) tanpa pemrosesan tambahan apapun untuk menguji kemampuan model CNN dalam mempelajari pola langsung dari piksel warna mentah.

\vspace{2mm}
\noindent\textbf{Komposit Softmap Spasial (\textit{Spatial Softmap Composite})} \\
Masukan tiga kanal kontinu yang dibentuk oleh: (1) \textit{softmap} pembuluh darah berbasis
plain Gabor, (2) \textit{softmap} area tepi (\textit{ridge softmap}) berbasis Meijering dan
\textit{oriented top-hat}, dan (3) kanal hijau ternormalisasi CLAHE yang telah di-masker FOV.
Semua kanal bersifat kontinu dan dimasker oleh FOV. Komposit ini dirancang sebagai bukti visual
terstruktur yang kaya bagi model CNN residual tanpa beban bias warna latar belakang retina.

Perbandingan visual dari kedua representasi citra input beserta kanal penyusun komposit softmap dapat dilihat pada Gambar \ref{fig:input_scenarios_comparison}.

\begin{figure}[H]
  \centering
  \includegraphics[width=\linewidth]{images/input_scenarios_comparison}
  \caption{Perbandingan visual representasi citra input beserta kanal pembentuk komposit softmap}
  \label{fig:input_scenarios_comparison}
\end{figure}

%--------------------------------------------------------------
\section{Baseline Klasik: Fitur Rekayasa Manual}
%--------------------------------------------------------------

Sebelum membangun model jaringan saraf, sebuah \textit{baseline} berbasis fitur rekayasa manual
(\textit{handcrafted features}) dibangun untuk menetapkan batas bawah performa yang harus
dilampaui. \textit{Baseline} ini menggunakan 48 fitur yang diekstrak dari setiap citra, yang mencakup:

\begin{itemize}
    \item \textbf{Fitur morfologi pembuluh darah}: Kerapatan pembuluh darah (\textit{vessel
    density}), proksi \textit{tortuositas}, anisotropi pembuluh, ukuran dan standar deviasi komponen.
    \item \textbf{Fitur tekstur}: Matriks kookurensi tingkat keabuan (\textit{Gray-Level
    Co-occurrence Matrix}, GLCM) \citep{Khandouzi2022} yang menghasilkan fitur kontras,
    korelasi, energi, dan homogenitas.
    \item \textbf{Fitur statistik intensitas}: Rata-rata dan standar deviasi intensitas pada
    setiap kanal warna.
    \item \textbf{Fitur area tepi (\textit{ridge})}: Luas, panjang, kerapatan, dan persentil
    respons dari komponen \textit{ridge} yang terdeteksi.
\end{itemize}

Fitur-fitur tersebut kemudian dimasukkan ke dalam tiga jenis model klasifikasi: \textit{Support
Vector Machine} (SVM) dengan kernel RBF, \textit{Random Forest} \citep{Breiman2001} dengan 400
pohon keputusan (\textit{n\_estimators=400, max\_depth=None}), dan \textit{Gradient Boosting}.
Evaluasi dilakukan menggunakan validasi silang 5-\textit{fold} yang distratifikasi berdasarkan kelas (tanpa pembatasan kelompok mata, berbeda dengan protokol \textit{group-aware} yang diterapkan pada model jaringan saraf).

%--------------------------------------------------------------
\section{Arsitektur Model Klasifikasi: TinyResNet}
%--------------------------------------------------------------

Model klasifikasi utama adalah arsitektur CNN residual kecil (\textit{TinyResNet}) yang
dirancang dari awal (\textit{from scratch}) tanpa menggunakan bobot pra-latih. Arsitektur ini
dipilih karena ukurannya yang proporsional dengan skala dataset sehingga risiko \textit{overfitting}
dapat dikendalikan.

Arsitektur \textit{TinyResNet} terdiri dari:
\begin{enumerate}
    \item Lapisan konvolusi awal $3 \times 3$ diikuti \textit{batch normalization} dan fungsi
    aktivasi ReLU.
    \item Tiga tahap \textit{residual block} bertingkat dengan lebar fitur masing-masing 48,
    96, dan 192 kanal. Setiap tahap menerapkan \textit{downsampling} spasial melalui konvolusi
    dengan \textit{stride} 2.
    \item \textit{Global Average Pooling} untuk meringkas peta fitur dua dimensi menjadi vektor
    satu dimensi.
    \item Lapisan \textit{Dropout} untuk regulasi.
    \item Lapisan \textit{Dense} dengan fungsi aktivasi \textit{softmax} untuk lima kelas.
\end{enumerate}

Struktur \textit{residual block} mengikuti pola dasarnya: \textit{conv-BN-ReLU-conv-BN} dengan
hubungan pintas (\textit{skip connection}) yang menjumlahkan masukan blok dengan keluarannya.
Arsitektur ini ditunjukkan pada Gambar \ref{fig:tinyresnet_arch}.

\begin{figure}[H]
  \centering
  \includegraphics[width=0.9\linewidth]{images/tinyresnet_arch}
  \caption{Struktur arsitektur model neural residual kecil (TinyResNet)}
  \label{fig:tinyresnet_arch}
\end{figure}

%--------------------------------------------------------------
\section{Protokol Pelatihan}
%--------------------------------------------------------------

\subsection{Fungsi Kerugian dan Pembobotan Kelas}

Pelatihan menggunakan \textit{focal loss} \citep{Lin2017} dengan parameter $\gamma = 1{,}0$
dan bobot kelas berbanding terbalik dengan frekuensi masing-masing kelas
(\textit{inverse-frequency class weights}). \textit{Focal loss} secara otomatis mengurangi
kontribusi sampel yang mudah diklasifikasikan, sehingga pelatihan lebih berfokus pada sampel
kelas minoritas yang sulit.

\subsection{Optimisasi dan Penjadwalan Laju Belajar}

Optimisasi menggunakan algoritma AdamW \citep{Loshchilov2019} dengan laju belajar awal
$10^{-3}$ dan peluruhan bobot (\textit{weight decay}) $5 \times 10^{-4}$. Laju belajar
dijadwalkan menggunakan strategi \textit{cosine annealing} selama 160 \textit{epoch}.

\subsection{Augmentasi Data}

Augmentasi data diterapkan selama pelatihan untuk meningkatkan variasi data latih dan
mengurangi \textit{overfitting}. Teknik augmentasi yang digunakan meliputi:
\begin{itemize}
    \item \textit{MixUp} \citep{Zhang2018} ($\alpha = 0{,}2$, probabilitas 0,5 per \textit{batch}):
    menggabungkan dua sampel citra dan labelnya secara proporsional menggunakan faktor pencampuran yang diambil dari distribusi Beta($\alpha$, $\alpha$).
    \item \textit{CutMix} \citep{Yun2019} ($\alpha = 1{,}0$, probabilitas 0,5 per \textit{batch}):
    menempelkan potongan gambar dari satu sampel ke sampel lain dengan label proporsional.
    \item \textit{Label smoothing} 0,1: meratakan distribusi target agar model tidak terlalu
    percaya diri (\textit{overconfident}) pada satu kelas.
    \item \textit{Stochastic depth} (\textit{DropPath}) \citep{Huang2016} dengan jadwal
    linear $0 \rightarrow 0{,}2$: menonaktifkan sebagian jalur residual secara acak
    selama pelatihan sebagai bentuk regulasi.
    \item \textit{Weight EMA} (\textit{Exponential Moving Average}) dengan peluruhan 0,999:
    merata-ratakan bobot model secara eksponensial untuk mendapatkan model yang lebih stabil.
    \item \textit{Flip test-time augmentation} pada saat inferensi OOF: prediksi dirata-ratakan
    dari citra asli dan cerminan horizontalnya.
\end{itemize}

\subsection{Protokol Validasi Silang}

Protokol utama yang digunakan adalah \textit{StratifiedGroupKFold} dengan 5 \textit{fold}.
Pembagian ini memastikan bahwa: (1) distribusi kelas pada setiap \textit{fold} kira-kira
proporsional, dan (2) seluruh citra dari kelompok yang sama (kelompok citra serupa yang
terdeteksi melalui NCC) selalu berada di sisi yang sama. Prediksi \textit{out-of-fold} (OOF, yaitu prediksi model pada data yang tidak diikutkan dalam pelatihan \textit{fold} tersebut) dari kelima \textit{fold} digabungkan untuk menghitung metrik evaluasi akhir.

%--------------------------------------------------------------
\section{Metrik Evaluasi}
%--------------------------------------------------------------

\subsection{Metrik Klasifikasi}

Metrik utama yang digunakan adalah \textit{macro F1-score}, yaitu rata-rata dari nilai
F1-score masing-masing kelas tanpa memperhitungkan jumlah sampel per kelas. Metrik ini dipilih
karena adanya ketidakseimbangan kelas pada dataset yang membuat akurasi keseluruhan menjadi
tidak cukup informatif. Kelas dengan sampel lebih banyak dapat mendominasi akurasi walaupun kelas
minoritas berkinerja buruk. Formula \textit{macro F1-score} adalah:

\begin{equation}
  \text{Macro F1} = \frac{1}{K} \sum_{k=1}^{K} F1_k = \frac{1}{K} \sum_{k=1}^{K}
  \frac{2 \cdot P_k \cdot R_k}{P_k + R_k}
  \label{eq:macro_f1}
\end{equation}

\noindent di mana $K$ adalah jumlah kelas, $P_k$ adalah presisi, dan $R_k$ adalah \textit{recall}
untuk kelas ke-$k$.

Selain \textit{macro F1-score}, metrik tambahan yang dilaporkan mencakup akurasi keseluruhan,
\textit{macro precision}, \textit{macro recall}, dan \textit{F1-score} per kelas.

